\usepackage{fontspec}
\usepackage{xeCJK}
\usepackage{unicode-math}
\usepackage{amsmath}
\usepackage{mathtools}
\usepackage{booktabs}
\usepackage{longtable}
\usepackage{array}
\usepackage{xcolor}
\usepackage{hyperref}

\defaultfontfeatures{Ligatures=TeX}

\setmainfont[
  Path=/overflow/htzhu/mingcheng_new/render_resources/chinese_math_pdf/fonts/texgyre-termes/,
  Extension=.otf,
  UprightFont=texgyretermes-regular,
  BoldFont=texgyretermes-bold,
  ItalicFont=texgyretermes-italic,
  BoldItalicFont=texgyretermes-bolditalic
]{texgyretermes-regular}

\setmathfont[
  Path=/overflow/htzhu/mingcheng_new/render_resources/chinese_math_pdf/fonts/texgyre-termes-math/,
  Extension=.otf
]{texgyretermes-math}

\setmathfont[
  Path=/overflow/htzhu/mingcheng_new/render_resources/chinese_math_pdf/fonts/newcomputermodern/,
  Extension=.otf,
  range={cal,bfcal}
]{NewCMMath-Regular}

\setCJKmainfont[
  Path=/overflow/htzhu/mingcheng_new/render_resources/chinese_math_pdf/texmf/fonts/opentype/public/noto-cjk/,
  Extension=.otf,
  UprightFont=NotoSerifSC-Regular,
  BoldFont=NotoSerifSC-Bold,
  AutoFakeSlant=0.18
]{NotoSerifSC-Regular}

\setCJKsansfont[
  Path=/overflow/htzhu/mingcheng_new/render_resources/chinese_math_pdf/texmf/fonts/opentype/public/noto-cjk/,
  Extension=.otf,
  UprightFont=NotoSansSC-Regular,
  BoldFont=NotoSansSC-Bold,
  AutoFakeSlant=0.18
]{NotoSansSC-Regular}

\setCJKmonofont[
  Path=/overflow/htzhu/mingcheng_new/render_resources/chinese_math_pdf/texmf/fonts/opentype/public/noto-cjk/,
  Extension=.otf,
  UprightFont=NotoSansSC-Regular,
  BoldFont=NotoSansSC-Bold
]{NotoSansSC-Regular}

\hypersetup{colorlinks=true,linkcolor=blue,urlcolor=blue,citecolor=blue}
\emergencystretch=3em
